% !TeX root = main.tex
\section{Torsion pairs and
decompositions of derived categories}
\label{sec:torsion-pairs-decomp}

We give criteria ensuring that a torsion pair in an abelian
category gives rise to a semi-orthogonal decomposition on the
level of derived categories. Our main assumption on the torsion
pair is that its two components are abelian subcategories of the
ambient abelian category. 
For bounded derived categories this is straightforward (see
Theorem~\ref{t:sod-DbBA}) as soon as some foundational results
for such torsion pairs are established (see
Proposition~\ref{p:abelian-torsion-pair}). Our result for
unbounded derived categories (see Theorem~\ref{t:sod-DBA}) needs
some technical assumptions on effectiveness of inverse and direct
truncation systems introduced in
Section~\ref{sec:notes-htpy-limits}. Fortunately, these
assumptions are satisfied for the derived category $\D_\qc(X)$ of
an algebraic stack (see Example~\ref{ex:ringed-topos-effective}
and Proposition~\ref{p:Dqc-truncation-complete}).


\subsection{Torsion pairs}
\label{sec:torsion-pairs}

We remind the reader of the notion of a \emph{torsion pair}. 
Sometimes
the terminology \emph{torsion theory} is used in the
literature. Standard references are
\cite[I.1]{beligiannis-reiten-torsion-theories},
\cite[1.12]{borceux-cat-2}, \cite{dickson-torsion}. 

\begin{definition}
  \label{d:torsion-pair}
  A \define{torsion pair} in an abelian category $\mathcal{B}$ is
  a pair 
  $(\mathcal{T}, \mathcal{F})$
  of strictly full subcategories of $\mathcal{B}$ such that the
  following two conditions hold:
  \begin{enumerate}[label=(TP{\arabic*})]
  \item 
    \label{enum:TP-ses}
    Any object $B \in \mathcal{B}$ fits into a short exact
    sequence
    \begin{equation*}
      0 \ra T \ra B \ra F \ra 0
    \end{equation*}
    with $T \in \mathcal{T}$ and $F \in \mathcal{F}$. 
  \item 
    \label{enum:TP-Hom-TF=0}
    We have $\Hom_\mathcal{B}(T, F)=0$ for all objects $T \in
    \mathcal{T}$ and $F \in \mathcal{F}$.
  \end{enumerate}
\end{definition}

\begin{remark}
  \label{rem:adjoints-f-and-t}
  Let $(\mathcal{T}, \mathcal{F})$
  be a torsion pair in an abelian category $\mathcal{B}$. 
  The axioms of a torsion pair immediately imply that 
  the short exact sequence in \ref{enum:TP-ses} is unique up to
  unique isomorphism and functorial in $B \in \mathcal{B}$.
  In particular, there are
  functors $t \colon \mathcal{B} \ra \mathcal{T}$ and $f \colon
  \mathcal{B} \ra \mathcal{F}$ together with 
  natural
  transformations
  $t \ra \id$ and $\id \ra f$ giving rise to short exact
  sequences 
  \begin{equation*}
    0 \ra t(B) \ra B \ra f(B) \ra 0,
  \end{equation*}
  for each object $B \in \mathcal{B}$. The object $B$ is in
  $\mathcal{T}$ (resp.\ $\mathcal{F})$ if and only if $f(B)=0$
  (resp.\ $t(B)=0$). 
  The functor $t$ is right adjoint to the inclusion functor
  $\mathcal{T} \ra \mathcal{B}$, and the functor $f$ is left
  adjoint to the inclusion functor 
  $\mathcal{F} \ra \mathcal{B}$.
\end{remark}

Recall that an \emph{abelian subcategory} of an abelian category
is a strictly full subcategory which is abelian and whose
inclusion functor is exact. 
A \emph{weak Serre subcategory}
is an abelian subcategory which is closed under
extensions in the ambient category.
If it is even closed under taking subobjects
and quotients in the ambient category, it is called a \emph{Serre
  subcategory}.  

\begin{remark}
  In the following we are only interested in torsion pairs 
  $(\mathcal{T}, \mathcal{F})$ where both
  $\mathcal{T}$ and $\mathcal{F}$ are abelian subcategories. 
  These torsion pairs are traditionally called
  \emph{hereditary} and \emph{cohereditary}. We do not use this
  terminology. 
\end{remark}

\begin{example}
  \label{e:A2-reps}
  Let $\mathcal{B}$ be the category of
  representations of the $A_2$-quiver $1 \ra 2$ over
  some ring. If $\mathcal{T}$ (resp.\ $\mathcal{F}$) denotes the
  subcategory of
  representations supported at the vertex $2$ (resp.\ $1$)
  then $(\mathcal{T}, \mathcal{F})$ is clearly a torsion pair in
  $\mathcal{B}$. Moreover, both $\mathcal{T}$ and
  $\mathcal{F}$ are abelian subcategories of $\mathcal{B}$.
\end{example}

\begin{proposition}
  \label{p:abelian-torsion-pair}
  Let $(\mathcal{T}, \mathcal{F})$
  be a torsion pair in an abelian category $\mathcal{B}$. If both
  $\mathcal{T}$ and $\mathcal{F}$ are abelian subcategories of
  $\mathcal{B}$, then the following statements are true.
  \begin{enumerate}
  \item 
    \label{enum:Hom-FT}
    % We have
    $\Hom_\mathcal{B}(F, T)=0$ for all objects $F \in
    \mathcal{F}$ and $T \in 
    \mathcal{T}$.
  \item 
    \label{enum:adjoints-exact}
    Both the right adjoint to the inclusion functor
    $\mathcal{T} \ra \mathcal{B}$ and 
    the left adjoint to the inclusion functor
    $\mathcal{F} \ra \mathcal{B}$ are exact,
    i.e., in the notation of Remark~\ref{rem:adjoints-f-and-t},
    both functors $t$ and $f$ are exact.
  \item 
    \label{enum:serre-subcats}
    Both $\mathcal{T}$ and $\mathcal{F}$ are Serre
    subcategories of $\mathcal{B}$.
  \item 
    \label{enum:Ext-TF}
    $\Ext^n_\mathcal{B}(T,F)=0$ for all objects
    $T \in \mathcal{T}$, $F \in \mathcal{F}$ and all integers $n
    \in \bZ$. 
  \end{enumerate}
\end{proposition}

\begin{remark}
  The most interesting part of
  Proposition~\ref{p:abelian-torsion-pair} 
  is certainly the $\Ext$-vanishing result in
  part~\ref{enum:Ext-TF}.
  After writing down its proof we learned that Brion has recently
  obtained the same result \cite[Lemma 2.3]{brion2018}.  
  For the convenience of the reader we give a full proof of
  Proposition~\ref{p:abelian-torsion-pair} even though parts
  \ref{enum:Hom-FT}, \ref{enum:adjoints-exact},
  \ref{enum:serre-subcats} are well-known or straightforward.
\end{remark}

\begin{proof}
  \ref{enum:Hom-FT}
  Let $\phi \colon F \ra T$ be a morphism from $F
  \in \mathcal{F}$ to $T \in \mathcal{T}$. 
  By \ref{enum:TP-ses}, we may
  put $\kernel(\phi)$ into a short exact sequence
  $0 \ra T' \ra \kernel(\phi) \ra F' \ra 0$ with $T' \in
  \mathcal{T}$ and $F' \in \mathcal{F}$. Since the composition of
  monomorphisms
  $T' \ra \kernel(\phi) \ra F$ is zero,
  by \ref{enum:TP-Hom-TF=0}, we have $T'=0$ and
  $\kernel(\phi) \in \mathcal{F}$.
  Similarly one shows $\cok(\phi) \in
  \mathcal{T}$. But then $\im(\phi)$ is the cokernel of
  $\kernel(\phi) \ra F$ and the
  kernel of $T \ra \cok(\phi)$ and hence is in $\mathcal{F} \cap
  \mathcal{T}$ since both $\mathcal{F}$ and $\mathcal{T}$ are
  abelian subcategories, i.e., $\im(\phi)=0$ and $\phi=0$.

  \ref{enum:adjoints-exact}
  We use the notation of
  Remark~\ref{rem:adjoints-f-and-t}. Clearly, $t$ and $f$ are
  additive. 
  Let $0 \ra B_1 \ra B_2 \ra B_3 \ra 0$ be a short exact
  sequence in $\mathcal{B}$. Then we obtain a commutative diagram
  with exact rows
  \begin{equation*}
    \xymatrix{
      {0} \ar[r]
      &
      {t(B_3)} \ar[r]
      &
      {B_3} \ar[r]
      &
      {f(B_3)} \ar[r]
      &
      {0}
      \\
      {0} \ar[r]
      &
      {t(B_2)} \ar[r] \ar[u]
      &
      {B_2} \ar[r] \ar[u]
      &
      {f(B_2)} \ar[r] \ar[u]
      &
      {0}
      \\
      {0} \ar[r]
      &
      {t(B_1)} \ar[r] \ar[u]
      &
      {B_1} \ar[r] \ar[u]
      &
      {f(B_1)} \ar[r] \ar[u]
      &
      {0.}
    }
  \end{equation*}
  We view the columns as complexes in $\mathcal{B}$ by adding
  zeroes.
  Then the diagram is a short exact sequence of
  complexes and gives rise to a long exact sequence on cohomology.
  Since the middle column is exact, since the cohomology objects
  of the left 
  (resp.\ right)
  column are in $\mathcal{T}$ (resp.\ $\mathcal{F})$, and since there
  are no nonzero morphisms 
  from $\mathcal{F}$ to $\mathcal{T}$, by \ref{enum:Hom-FT},
  the other two columns are also exact. This shows that $t$ and
  $f$ are exact.

  \ref{enum:serre-subcats}
  Let $F \ra B \ra F'$ be an exact sequence in $\mathcal{B}$ where
  $F$ and $F'$ are objects of 
  $\mathcal{F}$.
  To see that $\mathcal{F}$ is a Serre subcategory of
  $\mathcal{B}$ we need to see that $B \in \mathcal{F}$.
  The functor $t \colon \mathcal{B} \ra
  \mathcal{T}$ is exact, by 
  \ref{enum:adjoints-exact},
  and hence yields
  the exact sequence 
  $t(F) \ra t(B) \ra t(F')$ in $\mathcal{T}$.
  Since $t(F)=0=t(F')$ we obtain $t(B)=0$, i.e., $B \in
  \mathcal{F}$.

  The same argument proves that 
  $\mathcal{T}$ is a Serre subcategory of
  $\mathcal{B}$.

  \ref{enum:Ext-TF}
  The claim is trivially true for $n<0$. For 
  $n=0$ it is true 
  by axiom~\ref{enum:TP-Hom-TF=0}
  of a torsion pair.
  
  For $n=1$ we need to prove that 
  $\Ext^1_\mathcal{B}(T,F)=0$
  for $T \in \mathcal{T}$ and $F \in \mathcal{F}$. 
  Using Yoneda extensions (and
  \cite[\sptag{06XU}]{stacks-project})
  it is enough to show that any short exact sequence
  $0 \ra F \ra B \ra T \ra 0$ in $\mathcal{B}$ splits.
  Since $t$ is exact, by \ref{enum:adjoints-exact}, and $t(F)=0$,
  we obtain a morphism
  of short exact sequences
  \begin{equation*}
    \xymatrix{
      {0} \ar[r] 
      &
      {0} \ar[r] \ar[d]
      &
      {t(B)} \ar[r]^-{\sim} \ar[d]
      &
      {t(T)} \ar[r] \ar[d]^-{\sim}
      &
      {0}
      \\
      {0} \ar[r] 
      &
      {F} \ar[r]
      &
      {B} \ar[r]
      &
      {T} \ar[r]
      &
      {0}
    }
  \end{equation*}
  in $\mathcal{B}$; the vertical arrow on the right is an
  isomorphism since $T \in \mathcal{T}$ and hence $f(T)=0$.
  This shows that our sequence splits. (Since
  $\Hom_\mathcal{B}(T,F)=0$, the splitting is in fact unique.)

  Now let $n \geq 2$.
  Let $f \in \Ext_\mathcal{B}^n(T,F)$, i.e., $f \colon T \ra
  \Sigma^n F$ is a morphism in
  $\D(\mathcal{B})$.
  The $n$-extension $f$ can be written as the composition
  of a 1-extension and an $(n-1)$-extension: this is easy to
  prove directly or follows alternatively using
  Yoneda extensions and \cite[\sptag{06XU}]{stacks-project}.
  Anyhow, there is an object $B \in \mathcal{B}$ 
  such that $f$ is the composition
  \begin{equation*}
    f \colon T \xra{g} \Sigma^{n-1}B \xra{\Sigma^{n-1}h} \Sigma^n F,
  \end{equation*}
  where $g$ and $h \colon B \ra \Sigma F$ are 
  morphisms in $\D(\mathcal{B})$.
  The long exact $\Ext_\mathcal{B}(-,F)$-sequence obtained from
  the short 
  exact sequence $0 \ra t(B) \ra B \xra{\pi} f(B) \ra 0$ yields an
  isomorphism
  \begin{equation*}
    \pi^* \colon 
    \Ext^1_\mathcal{B}(f(B), F) \sira
    \Ext^1_\mathcal{B}(B, F) 
  \end{equation*}
  since we already know that 
  $\Ext^1_\mathcal{B}(t(B), F)$ and
  $\Hom_\mathcal{B}(t(B), F)$ vanish.
  Hence $h$ has the form $h=h' \circ \pi$ for a (unique) morphism
  $h' \colon f(B) \ra \Sigma F$ in $\D(\mathcal{B})$.
  Then
  \begin{equation*}
    f=(\Sigma^{n-1}h) \circ g = (\Sigma^{n-1}h') \circ
    (\Sigma^{n-1}\pi) \circ g. 
  \end{equation*}
  By induction, the $(n-1)$-extension $(\Sigma^{n-1}\pi) \circ g
  \colon T 
  \ra  \Sigma^{n-1}f(B)$ vanishes. Hence $f=0$.
\end{proof}

\subsection{Decompositions of bounded derived categories}
\label{sec:decomp-bound-deriv}

If $\mathcal{S}$ is a weak Serre subcategory of an abelian
category $\mathcal{B}$, the full subcategory of $\D(\mathcal{B})$
of objects with cohomology objects in
$\mathcal{S}$ is denoted
$\D_{\mathcal{S}}(\mathcal{B})$. 
Let 
$\D_{\mathcal{S}}^-(\mathcal{B})$,
$\D_{\mathcal{S}}^+(\mathcal{B})$,
and~$\D_{\mathcal{S}}^\bd(\mathcal{B})$ denote the full subcategories
of $\D_{\mathcal{S}}(\mathcal{B})$
of objects whose cohomology is bounded above, bounded below,
and bounded, respectively. 
All these categories are strictly full triangulated
subcategories of $\D(\mathcal{B})$, 
see, e.g., \cite[\sptag{06UQ}]{stacks-project}.
We refer the reader to \cite[Subsection~5.1]{BS-conservative}
for the notion of a semi-orthogonal decomposition.

\begin{theorem}
  \label{t:sod-DbBA}
  Let $\mathcal{B}$ be a weak Serre subcategory of an abelian
  category $\mathcal{A}$.
  Let 
  $\mathcal{T}$ and
  $\mathcal{F}$ be abelian subcategories of $\mathcal{B}$
  forming a torsion pair $(\mathcal{T}, \mathcal{F})$ in
  $\mathcal{B}$.
  Assume that 
  \begin{equation}
    \label{eq:DA-T-F}
    \Ext^n_{\mathcal{A}}(T, F)=0
  \end{equation}
  for all objects $T \in \mathcal{T}$, $F \in \mathcal{F}$ and
  all integers $n$. 
  Then 
  \begin{equation*}
    \D_{\mathcal{B}}^\bd(\mathcal{A})=
    \langle \D_{\mathcal{F}}^\bd(\mathcal{A}), 
    \D_{\mathcal{T}}^\bd(\mathcal{A})\rangle
  \end{equation*}
  is a semi-orthogonal decomposition.
\end{theorem}

\begin{remark}
Condition~\eqref{eq:DA-T-F} is automatically satisfied
in the special case when $\mathcal{B}=\mathcal{A}$,
by Proposition~\ref{p:abelian-torsion-pair}.\ref{enum:Ext-TF}.
In general, the condition is automatically satisfied for $n \leq 1$
by the same proposition, since $\mathcal{B}$ is closed under extensions.
\end{remark}

\begin{proof}
  By Proposition~\ref{p:abelian-torsion-pair}.\ref{enum:serre-subcats},
  $\mathcal{F}$ and $\mathcal{T}$ are Serre
  subcategories of $\mathcal{B}$, hence weak Serre
  subcategories of $\mathcal{A}$, and 
  $\D_{\mathcal{F}}^\bd(\mathcal{A})$ and  
  $\D_{\mathcal{T}}^\bd(\mathcal{A})$ are strictly full triangulated
  subcategories of $\D_\mathcal{B}^\bd(\mathcal{A})$.

  Repeated use of truncation shows that 
  $\D_{\mathcal{B}}^\bd(\mathcal{A})$ coincides with its smallest strictly
  full triangulated subcategory containing $\mathcal{B}$.
  Axiom~\ref{enum:TP-ses} shows that it is
  also its smallest such subcategory containing both
  $\mathcal{F}$ and $\mathcal{T}$.
  In particular, $\D_{\mathcal{B}}^\bd(\mathcal{A})$ is generated,
  as a triangulated category, by its subcategories
  $\D_{\mathcal{F}}^\bd(\mathcal{A})$ 
  and $\D_{\mathcal{T}}^\bd(\mathcal{A})$.

  By truncation and condition \eqref{eq:DA-T-F},
  there are no nonzero morphisms from 
  $\D_{\mathcal{T}}^\bd(\mathcal{A})$ to 
  $\D_{\mathcal{F}}^\bd(\mathcal{A})$.
\end{proof}
 
\subsection{Some notes on homotopy limits of truncations}
\label{sec:notes-htpy-limits}

Our aim is to state and prove Theorem~\ref{t:sod-DBA} 
in the following subsection~\ref{sec:decomp-unbo-deriv}
which is
the analog  
of Theorem~\ref{t:sod-DbBA} for unbounded
derived categories. In order to to this, 
we need some terminology and some basic observations.

We refer the reader to
\cite[Definition~1.6.4]{neeman-tricat} for the notion of a
homotopy colimit. The definition of a homotopy limit is dual.

Let $\mathcal{A}$ be an abelian category.
We propose the following terminology. 

\begin{definition}
\label{d:truncation-system}
\rule{0mm}{1mm}
\begin{enumerate}
\item
An inverse system 
\begin{equation*}
  (F_n)_{n \in \mathbb{N}} = 
  (\dots \ra F_{n+1} \ra F_n \ra \dots \ra F_2 \ra F_1 \ra F_0)
\end{equation*}
in $\D(\mathcal{A})$
is an \define{inverse truncation system} if
for each $n \in \mathbb{N}$
the object $F_n$ is in
$\D^{\geq -n}(\mathcal{A})$
and the map
$\tau_{\geq -n}F_{n+1} \to F_n$ induced by the transition
morphism is an isomorphism
(cf.\ Remark~\ref{rem:ITS-alternative} for a reformulation).
\item
If  $\D(\mathcal{A})$ has countable products, 
an inverse truncation system 
$(F_n)_{n \in \mathbb{N}}$ in $\D(\mathcal{A})$
is \define{effective} if
the induced maps $\HH^p(\holim_n F_n) \to
\lim_n\HH^p(F_n)$,
are isomorphisms for all $p \in \bZ$.
Our assumption ensures the existence of homotopy limits,
and all limits $\lim_n\HH^p(F_n)$
exist by 
Remark~\ref{rem:ITS-alternative} below.  
\end{enumerate}
Dually, we define a \emph{direct truncation system} in the
obvious way. 
If $\D(\mathcal{A})$ has countable coproducts we can 
talk about 
\emph{effectiveness} of direct truncation systems.
\end{definition}

\begin{example}
  \label{ex:ITS}
  Given any object $F$ in $\D(\mathcal{A})$, we get an inverse
  truncation system $(\tau_{\geq-n}F)_{n \in \mathbb{N}}$ whose
  transition morphisms are the natural maps between the truncations.
  Similarly, we get a direct truncation system 
  $(\tau_{\leq n}F)_{n \in \mathbb{N}}$.
\end{example}

\begin{definition}
  \label{d:standardITS}
  An inverse (resp.\ direct) truncation system is called
  \define{standard} if it is 
  isomorphic to an inverse (resp.\ direct) truncation system as
  in Example~\ref{ex:ITS}. 
\end{definition}

\begin{remark}
  Informally, one may think of an arbitrary
  inverse truncation system as obtained from a
  possibly non-existing object by truncation. 
  Effectiveness says that the natural candidate for such an
  object, the homotopy limit,
  indeed does the job. We explain this in
  Remark~\ref{rem:effective} below.
\end{remark}

\begin{remark}
  \label{rem:ITS-alternative}
  An inverse system $(F_n)_{n \in \mathbb{N}}$ in 
  $\D(\mathcal{A})$ is an inverse truncation system if and only
  if $\HH^p(F_n) = 0$ for all $p < -n$ and the maps
  $\HH^p(F_{n+1}) \ra \HH^p(F_n)$ induced by the
  transition maps are 
  isomorphisms for all $p \geq -n$.
  This just means that the induced inverse systems 
  $(\HH^p(F_{n}))_{n \in \mathbb{N}}$ look as follows.  
  \begin{equation*}
    \xymatrix@R0em{
      {\ldots} \ar[r]^-{\sim} 
      & {\HH^1(F_2)} \ar[r]^{\sim}
      & {\HH^1(F_1)} \ar[r]^{\sim}
      & {\HH^1(F_0)}
      \\
      {\ldots} \ar[r]^-{\sim} 
      & {\HH^0(F_2)} \ar[r]^{\sim}
      & {\HH^0(F_1)} \ar[r]^{\sim}
      & {\HH^0(F_0)}
      \\
      {\ldots} \ar[r]^-{\sim} 
      & {\HH^{-1}(F_2)} \ar[r]^{\sim}
      & {\HH^{-1}(F_1)} \ar[r]
      & {0}
      \\
      {\ldots} \ar[r]^-{\sim} 
      & {\HH^{-2}(F_2)} \ar[r]
      & {0} \ar[r]
      & {0}
     }
  \end{equation*}
  In particular, 
  then all limits 
  $\lim_n\HH^p(F_n)$ exist in $\mathcal{A}$.
\end{remark}

\begin{remark}
  \label{rem:effective}
  Assume that $\D(\mathcal{A})$ has countable products.
  Let $G:=\holim_n F_n$ be a homotopy limit of 
  an inverse truncation system $(F_n)_{n \in \bN}$.
  It comes together with morphisms
  $G \ra F_n$, for all $n \in \bZ$.
  Since $F_n \in \D^{\geq -n}(\mathcal{A})$ they come from unique
  morphisms  
  $\tau_{\geq -n}G \ra F_n$.
  Then, as a consequence
  of 
  Remark~\ref{rem:ITS-alternative},
  the following conditions are equivalent.
  \begin{enumerate}
  \item The inverse truncation system $(F_n)_{n \in \bN}$ is
    effective.
  \item 
    For all $p \in \bZ$ there is an integer $n \geq -p$ such that
    $\HH^p(G) \to
    \HH^p(F_n)$ is an isomorphism; this condition is then
    automatically 
    true for all 
    integers $n \geq -p$.
  \item The morphism
    $(\tau_{\geq -n} G)_{n \in \bN} \ra (F_n)_{n \in
      \bN}$ of inverse truncation systems is an isomorphism.
  \end{enumerate}
  In particular, any effective inverse truncation system is
  standard in the sense of Definition~\ref{d:standardITS}; more
  precisely, up to canonical isomorphism
  it is obtained by truncation from its homotopy limit.
\end{remark}

We address the natural question whether an object is the
homotopy limit of its truncations.  

\begin{lemma}[Objects as homotopy limits of their truncations]
  \label{l:holim-truncations}
  Let $F$ be an object of $\D(\mathcal{A})$ and assume that
  $\D(\mathcal{A})$ has countable products. 
  If the inverse truncation system $(\tau_{\geq -n}F)_{n \in \bN}$ is
  effective there is a triangle  
  \begin{equation*}
    F \ra \prod_{n \in \bN} \tau_{\geq -n} F \ra \prod_{n \in
      \bN} \tau_{\geq -n} F \ra \Sigma F
  \end{equation*}
  in $\D(\mathcal{A})$ which exhibits $F$ as a
  homotopy limit of its truncations and whose first morphism
  is induced by the truncation maps $F \ra \tau_{\geq -n}F$.
\end{lemma}

\begin{proof}
  Let $G$ be a homotopy limit 
  of the given inverse system
  with defining triangle
  \begin{equation*}
    G \ra \prod \tau_{\geq -n}F \ra \prod \tau_{\geq -n}F \ra
    \Sigma G.
  \end{equation*}
  The obvious map from $F$ to the second term of this triangle
  comes from a 
  (possibly non-unique) morphism $\gamma \colon F \ra G$ to its
  first term which makes the upper left triangle in the
  following diagram commutative.
  Since the map $G \ra \tau_{\geq -n}F$ factors through the
  truncation map $G \ra \tau_{\geq -n}G$, we obtain the
  commutative diagram
  \begin{equation*}
    \xymatrix{
      {F} \ar[r] \ar[d]^-{\gamma} & {\tau_{\geq -n}F}\\
      {G} \ar[r] \ar[ur] & {\tau_{\geq -n} G.} \ar[u]
    }
  \end{equation*}
  Its right vertical arrow is an isomorphism
  since the inverse truncation system $(\tau_{\geq -n}F)_{n \in
    \bN}$ is effective, by the equivalent statements in 
  Remark~\ref{rem:effective}.
  Hence the inverse of this arrow is the morphism 
  $\tau_{\geq -n}(\gamma)$ which is therefore an
  isomorphism. Since $n \in \bN$ was arbitrary, this implies that
  $\gamma$ is an isomorphism. The claim follows.
\end{proof}

We now provide instances of inverse and direct truncation
systems that are effective. 
The following proposition is essentially Remark~2.3 in
\cite{neeman-homotopy-limits}.

\begin{proposition}
\label{p:ab4-trunc}
  Let $\mathcal{A}$ be an abelian category. If $\mathcal{A}$ has
  countable products which are exact
  then any inverse truncation system in $\D(\mathcal{A})$ is
  effective.
  Dually, if $\mathcal{A}$ has exact countable coproducts
  then any direct truncation system in $\D(\mathcal{A})$ is
  effective.
\end{proposition}

\begin{proof}
  Since the statements are dual,
  it is enough to prove one of them.
  We prove the first statement.

  It is clear that $\D(\mathcal{A})$ has all countable products.
  Let $(F_n)_{n \in \mathbb{N}}$ be an inverse
  truncation system in $\D(\mathcal{A})$.
  Consider a homotopy limit $\holim_n F_n$ together with a
  defining triangle  
  \begin{equation*}
    \holim_n F_n \to \prod_n F_n \xrightarrow{1 - \sigma}
    \prod_n F_n \to \Sigma \holim_n F_n,
  \end{equation*}
  where $\sigma$ is the shift map.
  Since $\mathcal{A}$ has exact countable products,
  the functor $\HH^p(-)$ preserves countable products for all $p$.
  Hence, we get an exact sequence
  \begin{equation*}
    \dots
    \xrightarrow{\delta} 
    \HH^p(\holim_n F_n) \to 
    \prod_n \HH^p(F_n) \xrightarrow{1 - \sigma} 
    \prod_n \HH^p(F_n) \xrightarrow{\delta} 
    \HH^{p + 1}(\holim_n F_n)
    \ra \dots
  \end{equation*}
  in $\mathcal{A}$ where $\sigma$ is again the shift  map.
  Note that the inverse system $(\HH^p(F_n))_{n \in \mathbb{N}}$
  is zero for $n < -p$ and constant for $n \geq -p$ (cf.\
  Remark~\ref{rem:ITS-alternative}). 
  Hence Remark~\ref{rem:1-shift-split-epi} below shows that 
  all morphisms $1- \sigma$ in our exact sequence are
  epimorphisms, i.e., all connecting 
  morphisms $\delta$ 
  vanish. Hence the canonical morphism
  $$
  \HH^p(\holim_n F_n) \to \lim_n \HH^p(F_n)
  $$
  is an isomorphism, i.e., our inverse truncation system is effective.
\end{proof}

\begin{remark}
  \label{rem:1-shift-split-epi}
  If an inverse system $(G_n)_{n \in \bN}$ in an additive
  category with countable
  products is \emph{eventually constant} in the sense that
  there is an $N \in \bN$ 
  such that all
  transition morphisms $G_{n+1} \ra G_n$ for $n \geq N$ are
  isomorphisms, then the morphism ``identity minus shift''
  \begin{equation*}
    1-\sigma \colon \prod_{n \in \bN} G_n \ra \prod_{n \in \bN} G_n
  \end{equation*}
  is a split epimorphism: there is a morphism $f$ in the other
  direction such that $(1-\sigma) \circ f = \id$. We leave the
  easy proof of this fact to the reader.
\end{remark}

\begin{example}
  \label{ex:modules-effective}
  If $R$ is an arbitrary ring, then all 
  inverse and direct truncation systems in the derived category
  $\D(\Mod(R))$ of $R$-modules are effective, by 
  Proposition~\ref{p:ab4-trunc}.
\end{example}

\begin{example}
  \label{ex:ringed-topos-effective}
  If $(X, \mathcal{O})$ is a ringed topos,
  then the category of $\mathcal{O}$-modules is Grothendieck
  abelian. 
  In particular, it has arbitrary small products and coproducts and
  all coproducts are exact;
  in particular, any direct truncation system in $\D(X)$ is
  effective, by Proposition~\ref{p:ab4-trunc}. 
  Products, however, need not be exact.
  Nevertheless, certain inverse truncation systems may still be
  effective, see Proposition~\ref{p:Dqc-truncation-complete} below.
\end{example}

% If $X$ is an algebraic stack, we define the derived category
% $\D_\qc(X)$ as the full subcategory 
% of objects with quasi-coherent cohomology in the derived category
% of sheaves of $\mathcal{O}_X$-modules in the lisse--\'etale topos
% $X_\liset$. It is a 
% strictly full triangulated subcategory of this ambient derived
% category by \cite[\sptag{07B4}, \sptag{06UQ}]{stacks-project}.

The following proposition follows from the proof of
\cite[\sptag{0D6M}]{stacks-project},
which basically is an abstract version of a result by Bökstedt--Neeman~\cite[Lemma~5.3]{neeman-homotopy-limits}.

\begin{proposition}
  \label{p:Dqc-truncation-complete}
  Let $X$ be an algebraic stack.
  Then any inverse truncation system in $\D(X)$ with terms in
  $\D_\qc(X)$ is effective. 
\end{proposition}

\begin{proof}
  It is clear that $\D(X)$ has all products.
  Let $(F_n)_{n \in \mathbb{N}}$
  be an inverse truncation system 
  in $\D(X)$ whose terms $F_n$ lie in $\D_\qc(X)$.

  By the definition of an inverse truncation system, the 
  transition morphism
  $F_{n+1} \ra F_n$ factors as 
  $F_{n+1} \ra \tau_{\geq -n}F_{n+1} \sira F_n$, 
  and we obtain a triangle
  \begin{equation*}
    \Sigma^{n + 1}\HH^{-n-1}(F_{n+1})
    \to F_{n+1}
    \to F_n
    \to \Sigma^{n + 2}\HH^{-n-1}(F_{n+1})
  \end{equation*}
  for each $n \in \mathbb{N}$.
  Fix $p \in \bZ$. 
  Since the cohomology sheaves are assumed to be quasi-coherent,
  the functor $\HH^p(U, -)$ vanishes on the first and fourth term of
  this triangle 
  for 
  all $n \geq -p$ and all affine schemes $U$ which are smooth
  over $X$. 
  Hence $\HH^p(U, F_{n+1}) \ra \HH^p(U, F_n)$ is an isomorphism
  for all $n \geq -p$, i.e., 
  the inverse system 
  $(\HH^p(U, F_n))_{n \in \bN}$
  of abelian groups
  is eventually constant.
  By Remark~\ref{rem:1-shift-split-epi}, 
  the morphism ``identity minus shift''
  \begin{equation*}
    1-\sigma \colon \prod_n\HH^p(U,F_n)
    \ra \prod_n \HH^p(U,F_n)
  \end{equation*}
  is a (split) epimorphism. Its kernel is $\lim_n \HH^p(U,F_n)$.

  Note that the functors $\HH^q(U, -) = \HH^q(\Rd \Gamma(U, -|_U))$
  commute with 
  products since the functors restriction to $U$ and $\Rd
  \Gamma(U, -)$ are right 
  adjoints and products 
  are exact in the category of abelian groups.
  Hence applying $\HH^p(U, -)$
  to the defining triangle of the homotopy limit $\holim_n F_n$ 
  yields an exact sequence
  \begin{equation*}
    \dots
    \ra
    \HH^p(U, \holim_n F_n) \to 
    \prod_n \HH^p(U, F_n) \xrightarrow{1 - \sigma} 
    \prod_n \HH^p(U, F_n)
    \ra
    \dots
  \end{equation*}
  of abelian groups. Since all morphisms $1-\sigma$ are
  surjective, as observed above, this sequence splits into 
  short exact sequences, and
  we get isomorphisms
  \begin{equation*}
    \HH^p(U,\holim_n F_n) 
    \xrightarrow{\sim}
    \lim_n \HH^p(U,F_n)
    \xrightarrow{\sim}
    \HH^p(U,F_m)
  \end{equation*}
  for all $m \geq -p$.
  For an arbitrary object $G$ in $\D(X)$,
  the sheaf $\HH^p(G)$ is 
  the sheaf associated to the presheaf
  $U \mapsto \HH^p(U, G)$.
  Hence the canonical map $\HH^p(\holim_n F_n) \to
  \HH^p(F_m)$ 
  is an isomorphism for all $m \geq -p$.
  This means that our system is effective, 
  by Remark~\ref{rem:effective}.
\end{proof}

\subsection{Decompositions of unbounded derived categories}
\label{sec:decomp-unbo-deriv}

Now we can state and prove Theorem~\ref{t:sod-DBA}.

\begin{remark}
\label{rem:cohomology-split}
Let $\mathcal{B}$ be a weak Serre subcategory of an abelian
category $\mathcal{A}$
and let $\mathcal{T}$ and $\mathcal{F}$ be abelian subcategories of
$\mathcal{B}$ forming a torsion pair $(\mathcal{T}, \mathcal{F})$ in
$\mathcal{B}$.
Consider an arbitrary triangle
\begin{equation*}
  B' \to B \to B'' \to \Sigma B'
\end{equation*}
in $\D_\mathcal{B}(\mathcal{A})$.
Then we have $B'$ in $\D_\mathcal{T}(\mathcal{A})$
and $B''$ in $\D_\mathcal{F}(\mathcal{A})$
if and only if the long exact cohomology sequence splits up
into short exact sequences
\begin{equation*}
  0 \to \HH^p(B')
  \to \HH^p(B)
  \to \HH^p(B'')
  \to 0
\end{equation*}
with first term $\HH^p(B')$ in $\mathcal{T}$ and third term
$\HH^p(B'')$ in
$\mathcal{F}$, for each $p \in \bZ$. The ``if''-part is trivial,
and the ``only if''-part
is a simple consequence of the fact that $\Hom_\mathcal{B}(F, T)$
vanishes for each object $F$ in $\mathcal{F}$ and $T$
in $\mathcal{T}$, by 
Proposition~\ref{p:abelian-torsion-pair}.\ref{enum:Hom-FT}.
\end{remark}

\begin{theorem}
  \label{t:sod-DBA}
  Let $\mathcal{B}$ be a weak Serre subcategory of an abelian
  category $\mathcal{A}$.
  Let $(\mathcal{T}, \mathcal{F})$ be a torsion pair in
  $\mathcal{B}$ where both $\mathcal{T}$ and
  $\mathcal{F}$ are abelian subcategories of $\mathcal{B}$.
  Assume that 
  \begin{equation*}
    \Ext_\mathcal{A}^n(T, F)=0
  \end{equation*}
  for all objects $T \in \mathcal{T}$, $F \in \mathcal{F}$ and
  all integers $n \in \bZ$. 
  Assume that $\D(\mathcal{A})$ has countable products and
  coproducts. 
  If every inverse and every direct truncation
  system in 
  $\D(\mathcal{A})$ with terms in $\D_\mathcal{B}(\mathcal{A})$
  is effective, 
  in the sense of Definition~\ref{d:truncation-system},
  then 
  \begin{equation*}
    \D_{\mathcal{B}}(\mathcal{A})=
    \langle \D_{\mathcal{F}}(\mathcal{A}), 
    \D_{\mathcal{T}}(\mathcal{A})\rangle
  \end{equation*}
  is a semi-orthogonal decomposition.
\end{theorem}

\begin{proof}
\textbf{Semi-orthogonality}:
We claim that
\begin{equation}
  \label{eq:Hom-vanishing-unbounded}
  \Hom_{\D(\mathcal{A})}(L, R)=0
\end{equation}
for all objects $L \in \D_\mathcal{T}(\mathcal{A})$ and $R \in
\D_\mathcal{F}(\mathcal{A})$.

Our proof is a straightforward d\'evissage argument.
Assume first that $L$ is in $\D^-_\mathcal{T}(\mathcal{A})$ and that $R$ is
in $\D^+_\mathcal{F}(\mathcal{A})$.
Then any morphism $L \to R$ factors as $L \to \tau_{\geq n} L \to
\tau_{\leq m} R \to R$ for sufficiently small $n$ and
sufficiently large $m$.
Hence the morphism vanishes by Theorem~\ref{t:sod-DbBA}.

Keep the assumption on $L$ and assume that $R$ is an arbitrary
object of $\D_\mathcal{F}(\mathcal{A})$.
Effectiveness of inverse
truncation systems with terms in 
$\D_\mathcal{B}(\mathcal{A})$ and
Lemma~\ref{l:holim-truncations} provide a triangle
\begin{equation*}
  R \ra 
  \prod_{n \in \bN} \tau_{\geq -n} R 
  \ra
  \prod_{n \in \bN} \tau_{\geq -n} R 
  \ra
  \Sigma R
\end{equation*}
in $\D(\mathcal{A})$.
The functor $\Hom_{\D(\mathcal{A})}(L,-)$ vanishes on the second and
third term and their shifts by the universal property of
products, boundedness from below of
the truncations and the vanishing result already proven.
Hence it vanishes also on the first term as desired.

Similarly, 
effectiveness of direct
truncation systems with terms in 
$\D_\mathcal{B}(\mathcal{A})$
and the dual of
Lemma~\ref{l:holim-truncations}
give
\eqref{eq:Hom-vanishing-unbounded}
for arbitrary objects $L$ in $\D_\mathcal{T}(\mathcal{A})$ and
$R$ in $\D_\mathcal{F}(\mathcal{A})$.

\textbf{Completeness:}
To finish the proof, we need to show that any object $B$ in
$\D_{\mathcal{B}}(\mathcal{A})$ fits into a triangle
\begin{equation}
\label{eq:tria-complete}
T \to B \to F \to \Sigma T 
\end{equation}
with $T$ in $\D_{\mathcal{T}}(\mathcal{A})$ and
$F$ in $\D_{\mathcal{F}}(\mathcal{A})$.

Assume first that $B$ is an object in $\D^-_\mathcal{B}(\mathcal{A})$
and consider its associated standard inverse truncation system
$(B_n)_{n \in \mathbb{N}}$ where $B_n = \tau_{\geq -n} B$
(see Example~\ref{ex:ITS}).
Since the cohomology of each $B_n$ is bounded,
the semi-orthogonal decomposition in Theorem~\ref{t:sod-DbBA}
provides commutative squares
\begin{equation*}
  \xymatrix{
    T_n \ar[r]\ar[d] & T_{n+1}\ar[d] \\
    B_n \ar[r] & {B_{n+1},} 
  }
  \qquad
  \xymatrix{
    \HH^p(T_n) \ar[r]\ar[d] & \HH^p(T_{n+1})\ar[d] \\
    \HH^p(B_n) \ar[r] & \HH^p(B_{n+1}), 
  }
\end{equation*}
where each $T_n$ is the projection of $B_n$ to
$\D^\bd_\mathcal{T}(\mathcal{A})$ along this semi-orthogonal decomposition.
By Remark~\ref{rem:cohomology-split}, the vertical maps in the right
diagram are just the torsion subobjects with respect to the
torsion pair $(\mathcal{T}, \mathcal{F})$.
The characterization of inverse truncation systems in
Remark~\ref{rem:ITS-alternative} 
and the fact that 
$(B_n)_{n \in \mathbb{N}}$
is such a system make it 
clear that $(T_n)_{n \in \mathbb{N}}$
is such a system as well.
By assumption, it is effective, and hence obtained from its
homotopy limit $T$ in $\D(\mathcal{A})$ by truncation in a
canonical way, see
Remark~\ref{rem:effective}. 
Without loss of generality we may hence assume that 
$T_n=\tau_{\geq -n}T$ for all $n \in \bN$.
It is clear that $T \in \D^-_\mathcal{T}(\mathcal{A})$.
Note that $T$ and $B$ are the homotopy limits of 
the effective inverse truncation systems 
$(T_n)_{n \in \mathbb{N}}$ and $(B_n)_{n \in \mathbb{N}}$,
by Lemma~\ref{l:holim-truncations}.
Hence the morphism $(T_n)_{n \in \mathbb{N}} \ra (B_n)_{n \in \mathbb{N}}$
of inverse systems induces a morphism
$T \ra B$ which makes the following two squares commutative.
\begin{equation*}
  \xymatrix{
    {T} \ar[r]\ar[d] & {T_{n}} \ar[d] \\
    {B} \ar[r] & {B_{n}} 
  }
  \qquad
  \xymatrix{
    {\HH^p(T)} \ar[r]\ar[d] & {\HH^p(T_{n})}\ar[d] \\
    {\HH^p(B)} \ar[r] & {\HH^p(B_{n})} 
  }
\end{equation*}
Obviously, the horizontal arrows in the right
square are isomorphisms for all $n \geq -p$.
Hence $\HH^p(T) \to \HH^p(B)$ is the torsion
subobject with respect to the torsion pair $(\mathcal{T},
\mathcal{F})$, for each $p \in \bZ$.
Complete the morphism $T \to B$ to a triangle as in \eqref{eq:tria-complete}.
Then the associated long exact sequence on cohomology 
splits up into short exact sequences with $\HH^p(F)$ in $\mathcal{F}$.
In particular, the object $F$ is in $\D^-_\mathcal{F}(\mathcal{A})$.

This establishes the semi-orthogonal decomposition
$
\D^-_{\mathcal{B}}(\mathcal{A}) =
\langle \D^-_{\mathcal{F}}(\mathcal{A}), 
\D^-_{\mathcal{T}}(\mathcal{A})\rangle.
$
By the dual version of the argument above,
starting with an object $B$ in $\D_\mathcal{B}(\mathcal{A})$
and 
using the fact that all direct truncation systems with terms in
$\D_{\mathcal{B}}(\mathcal{A})$ are effective,
we get the desired semi-orthogonal decomposition of
$\D_\mathcal{B}(\mathcal{A})$. 
\end{proof}

%%% Local Variables: 
%%% mode: latex
%%% TeX-master: "main"
%%% End: 
